% Cuerpo editor-ready del nuevo capitulo 57.

\section{Défaut somme--produit d'APP et coefficient linéaire du produit modulaire d'Euler}
\label{sec:l9s102:euler-defecto}

\subsection{Identité exacte \texorpdfstring{\(459-405=54=[q]t_3(q)\)}{459-405=54=[q]t3(q)}}

Sur le support de \(81\) cellules, les deux feuillets d'APP produisent
\begin{equation}
 S_+=405,
 \qquad
 S_\times=459,
 \qquad
 S_\times-S_+=54.
 \label{eq:l9s102:euler-405-459}
\end{equation}
La régularisation triadique fournit l'exposant \(-1/12\). Sa
publication dodécaphasique détermine le pôle \(q^{-1}\), et le quotient de
Dedekind de niveau \(3\) s'écrit
\begin{equation}
 t_3(q)
 =\left(\frac{\eta(\tau)}{\eta(3\tau)}\right)^{12}
 =q^{-1}\prod_{n\ge1}(1+q^n+q^{2n})^{-12},
 \qquad q=e^{2\pi i\tau}.
 \label{eq:l9s102:euler-t3}
\end{equation}

\begin{theorem}[Identité somme--produit d'Euler--HMT]
\label{thm:l9s102:euler-54}
Le défaut entre les deux évaluations d'APP coïncide avec le coefficient
linéaire de \(t_3\) :
\begin{equation}
 \boxed{
 S_\times-S_+=54=[q]\,t_3(q).}
 \label{eq:l9s102:euler-54}
\end{equation}
\end{theorem}

\begin{proof}
Chaque facteur admet
\((1+q^n+q^{2n})^{-12}=1-12q^n+O(q^{2n})\). Pour le coefficient de
\(q\) postérieur au pôle \(q^{-1}\), les partitions de degré
\(2\) du produit de niveau \(3\) contribuent. Le développement fini certifié produit
\([q]t_3=54\). APP obtient le même entier par la soustraction exacte
\(459-405\).
\end{proof}

\subsection{Conservation du défaut au moyen du résidu, du quotient et de la retenue nonadique}

L'identité \cref{eq:l9s102:euler-54} relie la différence additive de deux
feuillets finis à un produit modulaire infini. Le lien conserve
l'origine de \(54\) : somme, produit et régularisation demeurent trois
coordonnées distinctes de l'état.

\subsubsection{Complétion EMRH et conservation de l'unité}
\label{sec:l9s102:euler-emrh}

La fibre ternaire de profondeur \(6\) possède \(3^6=729\) éléments. Sa
complétion décimale distingue
\begin{equation}
 729+270=999,
 \qquad
 729+271=1000,
 \qquad
 271=243+27+1.
 \label{eq:l9s102:euler-729-1000}
\end{equation}
Le terme terminal \(+1\) est l'unité de fermeture transmise par la
retenue. L'égalité binomiale
\begin{equation}
 1000=(9+1)^3=9^3+3\cdot9^2+3\cdot9+1
 \label{eq:l9s102:euler-binomial-1000}
\end{equation}
exprime le transport entre la carte nonadique et la publication décimale sans
perdre la généalogie de la complétion.

L'évaluation archimédienne est réalisée sur des cylindres semi-fermés
compatibles. Les deux développements d'un rationnel terminal représentent la
même valeur et conservent des histoires de bord distinctes. Cette structure
type l'identité \(0.999\ldots=1\) comme égalité d'évaluation accompagnée
d'une mémoire de voie.

\section{Algèbres quadratiques du TRIT et loi exponentielle elliptique, parabolique et hyperbolique}
\label{sec:l9s102:euler-algebras-trit}

\subsection{Famille opératorielle \texorpdfstring{\(J_{+1}^2=-I\), \(J_0^2=0\) et \(J_{-1}^2=I\)}{J+1^2=-I, J0^2=0 et J-1^2=I}}

Para \(\tau\in\{+1,0,-1\}\), definimos
\begin{equation}
 \mathbb A_\tau
 =\mathbb R[J_\tau]/(J_\tau^2+\tau I).
 \label{eq:l9s102:euler-algebras}
\end{equation}
La série exponentielle converge dans chaque algèbre et se réduit à
\begin{align}
 e^{tJ_{+1}}&=\cos t\,I+\sin t\,J_{+1},
 \label{eq:l9s102:euler-eliptica}\\
 e^{tJ_0}&=I+tJ_0,
 \label{eq:l9s102:euler-parabolica}\\
 e^{tJ_{-1}}&=\cosh t\,I+\sinh t\,J_{-1}.
 \label{eq:l9s102:euler-hiperbolica}
\end{align}
TRIT assigne respectivement fermeture elliptique, variation du premier ordre et
ouverture hyperbolique. L'identité classique occupe le feuillet \(\tau=+1\) :
\begin{equation}
 \boxed{\operatorname{Exp}(\pi J_{+1})+I=0.}
 \label{eq:l9s102:euler-clasica}
\end{equation}
Le feuillet parabolique conserve le premier jet, et le feuillet hyperbolique publie
un couple de valeurs propres réciproques.

\subsection{Identités \texorpdfstring{\(\operatorname{Exp}(\pi J_{+1})+I=0\)}{Exp(pi J+1)+I=0} et \texorpdfstring{\(\operatorname{Exp}(2\log\varphi\,H_\varphi)=F_{\mathrm{av}}^2\)}{Exp(2 log phi H_phi)=F_av^2}}

Soit
\begin{equation}
 F_{\mathrm{av}}=\begin{pmatrix}0&1\\1&1\end{pmatrix},
 \qquad
 H_\varphi=\frac{2F_{\mathrm{av}}^2-3I}{2\varphi-1}.
 \label{eq:l9s102:euler-hphi}
\end{equation}
Alors \(H_\varphi^2=I\) et
\begin{equation}
 \boxed{
 \operatorname{Exp}(2\log\varphi\,H_\varphi)=F_{\mathrm{av}}^2.}
 \label{eq:l9s102:euler-fav}
\end{equation}
Dans cette coordination, \(\pi\) mesure la fermeture angulaire, \(\varphi\) la
valeur propre expansive de la mémoire et \(e\) l'opération qui transforme un
générateur additif en transport multiplicatif.

\section{Principe d'Ambivalence : involution entre \texorpdfstring{\(\pi/\varphi^2\)}{pi/phi^2} et \texorpdfstring{\(\varphi/\pi^2\)}{phi/pi^2} et accumulation nonadique d'orientation}
\label{sec:l9s102:euler-ambivalencia}

\subsection{Action \texorpdfstring{\(\mathcal Q(x,y)=(x/y^2,y/x^2)\)}{Q(x,y)=(x/y^2,y/x^2)} sur les lecteurs surfacique et volumétrique}

Nous définissons le lecteur quadratique
\begin{equation}
 \mathcal Q(x,y)
 =\left(\frac{x}{y^2},\frac{y}{x^2}\right).
 \label{eq:l9s102:euler-q}
\end{equation}
Dans les coordonnées
\begin{equation}
 P=xy,
 \qquad O=\frac{x}{y},
 \label{eq:l9s102:euler-po}
\end{equation}
son action est
\begin{equation}
 P\circ\mathcal Q=P^{-1},
 \qquad
 O\circ\mathcal Q=O^3.
 \label{eq:l9s102:euler-q-action}
\end{equation}
La seconde itération satisfait
\begin{equation}
 \boxed{
 P\circ\mathcal Q^2=P,
 \qquad
 O\circ\mathcal Q^2=O^9.}
 \label{eq:l9s102:euler-q2}
\end{equation}
Le produit revient et l'orientation conserve une puissance nonadique de
mémoire.

\subsection{Identité \texorpdfstring{\(\varphi^3(\pi/\varphi^2)^2(\varphi/\pi^2)=1\)}{phi^3(pi/phi^2)^2(phi/pi^2)=1} et bidirectionnalité typée}

Pour \((x,y)=(\pi,\varphi)\), les deux orientations sont
\begin{equation}
 \rho_A=\frac{\pi}{\varphi^2},
 \qquad
 \rho_E=\frac{\varphi}{\pi^2},
 \label{eq:l9s102:euler-rho-a-e}
\end{equation}
et obéissent
\begin{equation}
 \boxed{\varphi^3\rho_A^2\rho_E=1.}
 \label{eq:l9s102:euler-ambivalencia-identidad}
\end{equation}
L'involution \((x,y)\mapsto(y,x)\) échange les orientations surfacique et
électronique du même lecteur \(R(x,y)=x/y^2\).

\section{Transformation exponentielle du générateur parangulaire \texorpdfstring{\(A_{\mathrm{rad}}I+C^*_{\mathrm{rad}}H_\varphi\)}{A_rad I+C*_rad H_phi}}
\label{sec:l9s102:euler-accion-par-angular}

La carte d'action
\begin{equation}
 T_{-34}(u)=u\,10^{-34}\mathcal U_S
 \label{eq:l9s102:euler-t34}
\end{equation}
transporte l'identité angulaire
\(C^*_{\deg}=2(\eta_{\mathrm{ret}}+\alpha_{\mathrm{HMT}})\) vers
\begin{equation}
 \boxed{
 T_{-34}(C^*_{\deg})
 =2\hbar_{\mathrm{ret}}
 +2T_{-34}(\alpha_{\angle,\deg}).}
 \label{eq:l9s102:euler-cstar-hbar}
\end{equation}
Chaque terme appartient ainsi à la même droite dimensionnelle
\(\mathcal L_S\).

\subsection{Projecteurs spectraux \texorpdfstring{\(P_\pm=(I\pm H_\varphi)/2\)}{P+/-=(I+/-H_phi)/2} et valeurs propres conjuguées \texorpdfstring{\(q_\pm\)}{q+/-}}

Soit
\begin{equation}
 B=A_{\mathrm{rad}}I+C^*_{\mathrm{rad}}H_\varphi,
 \qquad
 P_\pm=\frac{I\pm H_\varphi}{2}.
 \label{eq:l9s102:euler-b}
\end{equation}
L'exponentiation transforme le générateur additif en deux feuillets
multiplicatifs conjugués :
\begin{equation}
 \boxed{
 e^{-B}=q_+P_++q_-P_-,
 \qquad
 q_\pm=e^{-A_{\mathrm{rad}}\mp C^*_{\mathrm{rad}}}.}
 \label{eq:l9s102:euler-exp-b}
\end{equation}

\subsection{Compatibilité avec la publication parangulaire et les canaux \texorpdfstring{\(90/120\)}{90/120}}

Les cardinaux \(90\) et \(120\) proviennent de l'incidence tritique du TPK.
Le couple \(q_\pm\) ne les engendre pas : il évalue spectralement l'opérateur
\(V_{90,120}\) sur ces canaux préalablement construits.

\section{Bord elliptique d'action et intérieur hyperbolique de Hilbert--Beltrami--Klein}
\label{sec:l9s102:euler-elipse-klein}

\subsection{Coordonnées conjuguées de l'ellipse d'action et loi \texorpdfstring{\(\operatorname{Area}(\theta)=\hbar_{\mathrm{ret}}\theta\)}{Area(theta)=hbar_ret theta}}

La rapidité angulaire
\begin{equation}
 \eta=\operatorname{artanh}\!\left(\frac{C^*}{A}\right)
 \label{eq:l9s102:euler-rapidez}
\end{equation}
définit les demi-axes symplectiques
\begin{equation}
 a_S=\sqrt{2\hbar_{\mathrm{ret}}e^\eta},
 \qquad
 b_S=\sqrt{2\hbar_{\mathrm{ret}}e^{-\eta}}.
 \label{eq:l9s102:euler-semiejes}
\end{equation}
Avec
\begin{equation}
 Q=a_S\cos\theta,
 \qquad
 P=b_S\sin\theta,
 \label{eq:l9s102:euler-qp}
\end{equation}
l'action contenue dans l'arc satisfait
\begin{equation}
 \boxed{\operatorname{Area}(\theta)=\hbar_{\mathrm{ret}}\theta.}
 \label{eq:l9s102:euler-area-accion}
\end{equation}

\subsection{Correspondance \texorpdfstring{\(\rho=\tan(\theta/2)=\tanh(u/2)\)}{rho=tan(theta/2)=tanh(u/2)} entre bord périodique et intérieur hyperbolique}

La coordonnée de Gudermann
\begin{equation}
 \rho=\tan\frac\theta2=\tanh\frac u2
 \label{eq:l9s102:euler-gudermann}
\end{equation}
coordonne les cartes circulaire et hyperbolique.
Le bord porte l'orbite symplectique ; le domaine convexe intérieur porte
la métrique hyperbolique de Hilbert--Beltrami--Klein. La coordonnée de
Gudermann réalise leur correspondance radiale et conserve la distinction des
métriques.

\section{Compatibilité de la fermeture opératorielle d'Euler avec les publications archimédienne et exceptionnelle}
\label{sec:l9s102:euler-doble-proyeccion}

\subsection{Double projection du même état enrichi et conservation de l'incidence}

Soit \(X^{\mathrm{enr}}\) l'état enrichi construit par
APP--TRIT--TPK. Ses deux publications principales sont
\begin{equation}
 X^{\mathrm{enr}}
 \xrightarrow{\ \mathcal P_{\mathrm{arq}}\ }\mathbb R,
 \qquad
 X^{\mathrm{enr}}
 \xrightarrow{\ \mathcal P_{\mathrm{exc}}\ }
 \mathrm{Witt}\to\mathrm{Golay}\to\mathrm{Leech}
 \to V^\natural\to\mathbb M\to J.
 \label{eq:l9s102:euler-doble-proyeccion}
\end{equation}
La première contracte des cylindres compatibles vers la valeur archimédienne ; la
seconde conserve l'incidence de bord jusqu'au corridor exceptionnel.
L'entier \(54\) apparaît dans APP et dans le coefficient modulaire de
\cref{eq:l9s102:euler-54} ; le corridor \(4\to2\to6\to54\) transporte
ensuite cette incidence vers le réseau de Leech, l'algèbre d'opérateurs
de vertex et Moonshine. L'action du Monstre réside dans
\(V^\natural\) ; la généalogie des chiffres et des voies demeure dans le domaine
du TPK et atteint le corridor au moyen des applications d'incidence déclarées.

\subsection{Système macrovectoriel \texorpdfstring{\(\mathfrak E_{\mathrm{HMT}}=0\)}{E_HMT=0} : confluence des identités eulériennes dans un unique opérateur}
\label{sec:l9s102:euler-macrovector}

Les identités précédentes sont réunies dans l'application typée
\begin{equation}
 \boxed{
 \mathfrak E_{\mathrm{HMT}}=
 \begin{pmatrix}
 (S_\times-S_+)-[q]t_3\\[1mm]
 \operatorname{Exp}(I)-eI\\
 \operatorname{Exp}(\pi J_{+1})+I\\
 \operatorname{Exp}(J_0)-I-J_0\\
 \operatorname{Exp}(2\log\varphi\,H_\varphi)-F_{\mathrm{av}}^2\\
 \varphi^3\rho_A^2\rho_E-1\\
 T_{-34}(C^*)-2\hbar_{\mathrm{ret}}
 -2T_{-34}(\alpha_{\angle,\deg})
 \end{pmatrix}=0.}
 \label{eq:l9s102:euler-macrovector}
\end{equation}
Chaque composante conserve son domaine et sa propre démonstration. L'application
agrégée met en évidence la coordination de la somme, du produit, de la fermeture elliptique, du premier
jet, de l'expansion hyperbolique, de l'ambivalence et de l'action.

La composition interne APP--TRIT--TPK transporte déjà somme, produit,
orientation, résidu, quotient, retenue, voie et mémoire jusqu'aux
publications précédentes. On ne postule donc pas d'opérateur d'entrelacement HMT
absent. Une promotion ultérieure ne peut être formulée qu'au moyen du carré
typé spécifique exigé par une représentation ultérieure ---par exemple,
une équivariance complète de groupe jusqu'à Witt---, sans amoindrir la confluence
opératorielle démontrée ici.
